\documentclass[11pt]{article}
\usepackage[a4paper,margin=21mm]{geometry}
\usepackage[no-math]{fontspec}
\setmainfont{Latin Modern Roman}
\newfontfamily\CJKfont{Noto Serif CJK SC}
\XeTeXlinebreaklocale "zh"
\XeTeXlinebreakskip=0pt plus 1pt
\usepackage{amsmath,amssymb,amsthm,amscd,graphicx,color}
\usepackage{xurl}
\usepackage[unicode,hidelinks]{hyperref}
\allowdisplaybreaks
\setlength\emergencystretch{3em}
\numberwithin{equation}{section}

\numberwithin{equation}{section}

\newtheorem{Theorem}{Theorem}[section]
\newtheorem{Corollary}[Theorem]{Corollary}
\newtheorem{Lemma}[Theorem]{Lemma}
\newtheorem{Proposition}[Theorem]{Proposition}
 { \theoremstyle{definition}
\newtheorem{Definition}[Theorem]{Definition}
\newtheorem{Note}[Theorem]{Note}
\newtheorem{Example}[Theorem]{Example}
\newtheorem{Remark}[Theorem]{Remark} }

\newcommand\re[1]{(\ref{#1})}
\newcommand{\op}{\oplus}
\newcommand{\ot}{\otimes}

\renewcommand{\L}{\mathbf{L}}
\renewcommand{\P}{\mathbf{P}}

\def\Gbdle{\mathcal{G}}
\def\Hbdle{\mathcal{H}}
\def\isom{\cong}
\def\galg{\mathfrak{g}}
\def\falg{\mathfrak{f}}
\def\balg{\mathfrak{b}}
\def\palg{\mathfrak{p}}
\def\halg{\mathfrak{h}}
\def\malg{\mathfrak{m}}
\def\nalg{\mathfrak{n}}
\def\salg{\mathfrak{s}}
\def\som{\mathfrak{so}(\mathfrak{m})}
\def\sohn{\mathfrak{so}(\widehat{\mathfrak{n}})}
\def\dsum{\oplus}

\DeclareMathOperator{\Conf}{Conf}
\DeclareMathOperator{\Stab}{Stab}
\DeclareMathOperator{\Iso}{\mathrm{Isom}}
\DeclareMathOperator{\Tran}{\mathrm{Tran}}
\DeclareMathOperator{\id}{id}

\DeclareMathOperator{\E}{\mathbf{Euc}}
\DeclareMathOperator{\Hei}{\mathbf{Hei}}

\newcommand{\marg}[1]{\marginpar{\small {\tt #1 }}}
\newcommand{\eeqref}[1]{equation~(\ref{#1})}
\newcommand{\der}{{\mathrm d}}
\DeclareMathOperator{\hess}{\mathrm{Hess}}
\newcommand{\p}{\mathfrak{p}}
\newcommand{\Q}{\ensuremath{\mathbb{H}}}
\newcommand{\N}{\ensuremath{\mathbb{N}}}
\newcommand{\Z}{\ensuremath{\mathbb{Z}}}
%\newcommand{\R}{\ensuremath{\mathbb{R}}}
\newcommand{\C}{\ensuremath{\mathbb{C}}}
\newcommand{\K}{\ensuremath{\mathbb{K}}}
\renewcommand{\O}{\ensuremath{\mathrm{O}}}
%\newcommand{\R}{\ensuremath{\mathbb{R}}}
\renewcommand{\SS}{\ensuremath{\mathbb{S}}}
\newcommand{\vect}[1]{\boldsymbol{#1}}
\newcommand{\PSU}{\mathrm{PSU}}
\newcommand{\PSp}{\mathrm{PSp}}

\newcommand{\bmat}{\begin{pmatrix}}
\newcommand{\emat}{\end{pmatrix}}
\newcommand{\SSO}{\mathrm{SO}}
\newcommand{\CO}{\mathrm{CO}}
\newcommand{\SSp}{\mathrm{Sp}}
\newcommand{\SSU}{\mathrm{SU}}
\newcommand{\U}{\mathrm{U}}
\newcommand{\SSL}{\mathrm{SL}}
\newcommand{\GL}{\mathrm{GL}}
\newcommand{\1}{\mathbf{1}}
\newcommand{\I}{\mathbf{I}}
\newcommand{\ipq}{\1_{p,q}}
\newcommand{\J}{\mathbb{J}}
\newcommand{\heta}{\hat{\theta}}
\newcommand{\ups}{\Upsilon}
\newcommand{\e}{\mathrm{e}}
\newcommand{\ev}{\mathbf{e}}
\renewcommand{\d}{{\mathrm d}}
\newcommand{\bcase}{\begin{case}}
\newcommand{\ecase}{\end{case}}
\newcommand{\setcase}{\setcounter{case}{0}}
\newcommand{\bclaim}{\begin{claim}}
\newcommand{\eclaim}{\end{claim}}
\newcommand{\setclaim}{\setcounter{claim}{0}}
\newcommand{\bstep}{\begin{step}}
\newcommand{\estep}{\end{step}}
\newcommand{\setstep}{\setcounter{step}{0}}
\newcommand{\bhlem}{\begin{hlem}}
\newcommand{\ehlem}{\end{hlem}}
\newcommand{\sethlem}{\setcounter{hlem}{0}}

\newcommand{\Mat}{\mathrm{Mat}}

\newcommand{\bleer}{\begin{leer}}
\newcommand{\eleer}{\end{leer}}
\newcommand{\ul}{\underline}
\newcommand{\ol}{\overline}
\newcommand{\tbf}{\textbf}
\newcommand{\mc}{\mathcal}
\newcommand{\mb}{\mathbb}
\newcommand{\mf}{\mathfrak}
\newcommand{\bs}{\begin{Proposition}}
\newcommand{\es}{\end{Proposition}}
\newcommand{\btheo}{\begin{Theorem}}
\newcommand{\etheo}{\end{Theorem}}
\newcommand{\bfolg}{\begin{Corollary}}
\newcommand{\efolg}{\end{Corollary}}
\newcommand{\blem}{\begin{Lemma}}
\newcommand{\elem}{\end{Lemma}}
\newcommand{\bnote}{\begin{note}}
\newcommand{\enote}{\end{note}}
\newcommand{\bprf}{\begin{proof}}
\newcommand{\eprf}{\end{proof}}
\newcommand{\bd}{\begin{displaymath}}
\newcommand{\ed}{\end{displaymath}}
\newcommand{\be}{\begin{eqnarray*}}
\newcommand{\ee}{\end{eqnarray*}}
\newcommand{\eeqa}{\end{eqnarray}}
\newcommand{\beqa}{\begin{eqnarray}}
\newcommand{\bi}{\begin{itemize}}
\newcommand{\ei}{\end{itemize}}
\newcommand{\bnum}{\begin{enumerate}}
\newcommand{\enum}{\end{enumerate}}
\newcommand{\la}{\langle}
\newcommand{\ra}{\rangle}
\newcommand{\eps}{\epsilon}
\newcommand{\ve}{\varepsilon}
\newcommand{\vp}{\varphi}
\newcommand{\lra}{\longrightarrow}
\newcommand{\Lra}{\Leftrightarrow}
\newcommand{\Ra}{\Rightarrow}
\newcommand{\sub}{\subset}
\newcommand{\ems}{\emptyset}
\newcommand{\sms}{\setminus}
\newcommand{\ints}{\int\limits}
\newcommand{\sums}{\sum\limits}
\newcommand{\lims}{\lim\limits}
\newcommand{\bcup}{\bigcup\limits}
\newcommand{\bcap}{\bigcap\limits}
\newcommand{\beq}{\begin{equation}}
\newcommand{\eeq}{\end{equation}}
\newcommand{\einhalb}{\frac{1}{2}}
\newcommand{\rr}{\mathbb{R}}
\newcommand{\rn}{\mathbb{R}^n}
\newcommand{\ccc}{\mathbb{C}}
\newcommand{\cn}{\mathbb{C}^n}
\newcommand{\M}{M}
\newcommand{\drehgleich}{\mbox{\begin{rotate}{90}$=$ \end{rotate}}}
\newcommand{\turngleich}{\mbox{\begin{turn}{90}$=$ \end{turn}}}
\newcommand{\turnsimeq}{\mbox{\begin{turn}{270}$\simeq$ \end{turn}}}
\newcommand{\vf}{\varphi}
\newcommand{\earr}{\end{array}\]}
\newcommand{\barr}{\[\begin{array}}
\newcommand{\bvec}{\left(\begin{array}{c}}
\newcommand{\evec}{\end{array}\right)}
\newcommand{\sumk}{\sum_{k=1}^n}
\newcommand{\sumi}{\sum_{i=1}^n}
\newcommand{\suml}{\sum_{l=1}^n}
\newcommand{\sumj}{\sum_{j=1}^n}
\newcommand{\sumij}{\sum_{i,j=1}^n}
\newcommand{\sumkl}{\sum_{k,l=1}^n}
\newcommand{\suminf}{\sum_{k=0}^\infty}
\newcommand{\inv}{\frac{1}}
\newcommand{\g}{\mathfrak{g}}
\renewcommand{\a}{\mathfrak{a}}
\newcommand{\he}{\mathfrak{he}}
\newcommand{\siml}{\mathfrak{sim}}
\newcommand{\euc}{\mathfrak{euc}}
\renewcommand{\l}{\mathfrak{l}}
\newcommand{\h}{\mathfrak{h}}
\newcommand{\z}{\mathfrak{z}}
\newcommand{\m}{\mathfrak{m}}
\newcommand{\n}{\mathfrak{n}}
\renewcommand{\k}{\mathfrak{k}}
\renewcommand{\sl}{\mathfrak{sl}}
\newcommand{\hol}{\mathfrak{hol}}
\newcommand{\stab}{\mathfrak{stab}}
\newcommand{\Hol}{\mathrm{Hol}}
\newcommand{\+}{\oplus}
\renewcommand{\r}{\vect{r}}
\newcommand{\oo}{\vect{0}}
\newcommand{\x}{\vect{x}}
\newcommand{\lx}{\ltimes}
\newcommand{\rrn}{\mathbb{R}^n}
\newcommand{\so}{\mathfrak{so}}
\newcommand{\spin}{\mathfrak{spin}}
\newcommand{\co}{\mathfrak{co}}
\renewcommand{\b}{\flat}
\newcommand{\son}{{\mathfrak{so}(n)}}
\newcommand{\gl}{\mathfrak{gl}}
\newcommand{\lasl}{\mathfrak{sl}}
\renewcommand{\sp}{\mathfrak{sp}}
\newcommand{\su}{\mathfrak{su}}
\newcommand{\SU}{\mathbf{SU}}
\renewcommand{\U}{\mathbf{U}}
\newcommand{\SO}{\mathbf{SO}}
\newcommand{\SL}{\mathbf{SL}}
%\renewcommand{\u}{\mathfrak{u}}
\newcommand{\w}{\omega}
\newcommand{\pmh}{{\cal P}(M,h)}
\newcommand{\s}{\sigma}
\newcommand{\del}{\partial}
\newcommand{\ddt}{\frac{\partial}{\partial t}}
\newcommand{\ddr}{\frac{\partial}{\partial r}}
\newcommand{\ddx}{\frac{\partial}{\partial x}}
\newcommand{\ddz}{\frac{\partial}{\partial z}}
\newcommand{\ddu}{\frac{\partial}{\partial u}}
\newcommand{\ddv}{\frac{\partial}{\partial v}}
\newcommand{\du}{\partial_u}
\newcommand{\dv}{\partial_v}
\newcommand{\bbR}{\mathbb{R}}

\newcommand{\ddi}{\frac{\partial}{\partial y_k}}
\newcommand{\ddj}{\frac{\partial}{\partial y_l}}
\newcommand{\ddk}{\frac{\partial}{\partial y_k}}
\newcommand{\ddp}{\frac{\partial}{\partial p_k}}
\newcommand{\ddq}{\frac{\partial}{\partial q_k}}
\newcommand{\xz}{^{(x,z)}}
\newcommand{\mh}{(M,h)}
\newcommand{\wxz}{W_{(x,z)}}
\newcommand{\qmh}{{\cal Q}(M,h)}
\newcommand{\bbem}{\begin{Remark}}
\newcommand{\ebem}{\end{Remark}}
\newcommand{\bbez}{\begin{bez}}
\newcommand{\ebez}{\end{bez}}
\newcommand{\bbsp}{\begin{bsp}}
\newcommand{\ebsp}{\end{bsp}}
\newcommand{\pr}{\mathrm{pr}}
\newcommand{\widehats}{\hat{\s}}
\newcommand{\whut}{\w^{\widehats}}
\newcommand{\bhg}{{\cal B}_H(\lag)}
\newcommand{\aaa}{\rho}
\newcommand{\bb}{\s}
\newcommand{\laa}{\mf{a}}
%newcommand{\rho}{\rho}
\newcommand{\lam}{\lambda}
\newcommand{\LL}{\Lambda}
\newcommand{\D}{\Delta}
%\newcommand{\ß}{\s}
%\newcommand{\ä}{\rho}
\newcommand{\W}{\Omega}
\newcommand{\esel}{\ensuremath{\mathfrak{sl}(2,\ccc)}}
\newcommand{\kg}{{\cal K}(\lag)}
%\newcommand{\bg}{{\cal B}_h(\lag)}
\newcommand{\kk}{ \mathbb{K}}
\newcommand{\xy}{[x,y]}
\newcommand{\perdef}{$\stackrel{\text{\tiny def}}{\iff}$}
\newcommand{\eqdef}{\stackrel{\text{\tiny def}}{=}}
\newcommand{\lai}{\mf{i}}
\newcommand{\lar}{\mf{r}}
\newcommand{\Dim}{\mathsf{dim\ }}
\newcommand{\im}{\mathrm{im}}
\newcommand{\Ker}{\mathrm{Ker }}
\newcommand{\pro}{\text{pr}}
\newcommand{\trace}{\mathrm{tr }}
\DeclareMathOperator{\grad}{\mathrm{grad}}
\newcommand{\wt}{\widetilde}
\newcommand{\tnab}{\wt{\nabla}}
\newcommand{\nabg}{\nabla}
\newcommand{\nabh}{\nabla^h}
\newcommand{\nabk}{\nabla^{\cal K}}
\newcommand{\tM}{\widetilde{M}}
\newcommand{\hm}{\widehat{\M}}
\newcommand{\hnab}{\widehat{\nabla}}
\newcommand{\hR}{\widehat{\RR}}


\newcommand{\ten}{\widetilde{\cal N}}
\newcommand{\tg}{\widetilde{g}}
\newcommand{\tgg}{\widetilde{\g}}
\newcommand{\nabt}{\ol{\nabla}}
\newcommand{\ro}{\mathsf{P}}
\newcommand{\hro}{\hat{\mathsf{P}}}
\newcommand{\gtwo}{\lag_{2(2)}}
\renewcommand{\=}{\equiv}

\DeclareMathOperator{\Span}{\mathrm{span}}

\newcommand{\Ad}{\mathrm{Ad}}
\newcommand{\ad}{\mathrm{ad}}
\newcommand{\G}{\mathrm{G}}


\newcommand{\hook}{\makebox[7pt]{\rule{6pt}{.3pt}\rule{.3pt}{5pt}}\,}

\newcommand{\lecturecount}{\begin{center}{\sf (Lecture \Roman{lecturenr})}\end{center}\addtocontents{toc}{{\sf (Lecture \Roman{lecturenr})}}~\refstepcounter{lecturenr}
}
%% CHOL
\newcommand{\T}{\mathbb{T}}
\newcommand{\cur}{{\cal R}}
\newcommand{\pd}{{\cal P}}
%\newcommand{\Ro}{\mathsf{P}}


\renewcommand{\gg}{g}
\newcommand{\hg}{\widehat{\gg}}
\newcommand{\hh}{\mathrm{h}}
\newcommand{\II}{\mathrm{II}}
%\newcommand{\W}{\mathrm{W}}
\newcommand{\scal}{\mathrm{scal}}
\renewcommand{\t}{{(t)}}

%\newcommand{\Hess}{\mathrm{Hess}}

\newcommand{\bM}{{\overline{\cal M}}}
\newcommand{\bg}{{\overline{\mathbf{g}}}}
\newcommand{\bnab}{{\overline{\nabla}}}
\newcommand{\bR}{{\overline{\mathrm{R}}}}
\newcommand{\RR}{{R}}
\newcommand{\R}{\mathbb{R}}
\newcommand{\hRic}{\widehat{Ric}}
\newcommand{\bscal}{{\overline{\mathrm{scal}}}}

%\renewcommand{\epsilon}{\varepsilon}

\newcommand{\tcom}[2]{
\underline{#1} \marginpar{\small {\tt #2 }}}


\newcommand{\Id}{\mathrm{Id}}

\newcommand{\hei}{\mathfrak{hei}}

\newcommand{\inter}{\makebox[7pt]{\rule{6pt}{.3pt}\rule{.3pt}{5pt}}\,}
\newcommand{\dna}{d^\nabla\hspace{-3pt} A}
\newcommand{\ts}{\widetilde{\cal S}}
\renewcommand{\div}{\mathrm{div}}
%\renewcommand{\i}{\mathrm{i}}
%\newcommand{\insrt}{\mbox{\begin{turn}{180} $\invneg$ \end{turn}}}
%\newcommand{\turngleich}{\mbox{\begin{turn}{90}$=$ \end{turn}}}
%\newcommand{\turnsimeq}{\mbox{\begin{turn}{270}$\simeq$ \end{turn}}}
\newcommand{\liealg}[1]{\mathfrak{#1}}
\newcommand{\belabel}[1]{\begin{equation}\label{#1}}

\begin{document}
\begin{center}\Large Compact locally conformally Ricci-flat pseudo-Kähler examples on S1 times S(8n-1) with essential conformal transformations in signature (4n,4n)\end{center}
\noindent\textbf{Stable ID:} 1049\quad\textbf{Author:} Vicente Cortés; Thomas Leistner\par
\noindent\textbf{Work:} Compact Locally Conformally Pseudo-Kähler Manifolds with Essential Conformal Transformations\par
\noindent\textbf{Source:} \url{https://sigma-journal.com/2024/084/}\par
\noindent\textbf{Version:} Official SIGMA 20 (2024) article 084, DOI 10.3842/SIGMA.2024.084; journal PDF and native TeX acquired 2026-09-30, exact bytes pinned by SHA256\par
\noindent\textbf{Rights:} CC BY 4.0. Authors retain copyright. \url{https://creativecommons.org/licenses/by/4.0/}. Reuse and typesetting adaptation; original language and mathematical content preserved.\par
\noindent\textbf{Difficulty (prior selection preserved):} {\CJKfont H2: The author builds an indefinite Kähler metric with constant curvature components and verifies Ricci flatness and nonflatness explicitly. An anisotropic holomorphic homothety gives a compact cyclic quotient, while a commuting one-parameter homothety with fixed points proves essentiality in the induced conformal class. The metric construction and quotient/essentiality argument together form the higher-level obligation.}\par
\medskip\noindent\textbf{Editorial boundary note (not author text):} Complete original Theorem2.5 for non-zero symmetric sigma and all n,a,b hypotheses: the compact cyclic quotient, product topology, locally conformally Ricci-flat pseudo-Kähler conformal class of neutral signature, integrable complex structure and essential conformal transformations. The complete current metric construction, components/inverse, Christoffel/curvature tables, Ricci-flatness and locally/globally symmetric-space proofs, geodesic completeness and full nondegenerate holonomy argument are retained. The exact anisotropic holomorphic homotheties and their factorization proof, spherical–ellipsoidal fundamental domain, fixed-point essentiality and quotient lifting contradiction are included. Nontrivial fixed-point homothety essentiality and standard Kähler connection/Ricci identities remain precisely cited author prerequisites. One main compact conformal construction, explicitly announced as the main result; the transvection-algebra proposition, real-signature constructions and further product extensions are excluded and not counted. Literal author source notation and apparent index/dimension quirks remain unchanged, with counter reset only across the excluded transvection proposition. No proof generation, formula repair or independent mathematical review.\par
\medskip\hrule\medskip\noindent\textbf{Author text begins below.}\par
\setcounter{section}{1}
% Original source lines 390--398
A {\em conformal diffeomorphism} between semi-Riemannian manifolds $(M,g)$ and $(N,h)$ is a diffeomorphism
$\phi\colon M\to N$ such that
\[\phi^*h=\e^{2\vf}g,\]
for a smooth function $\vf\in C^\infty(M)$. A {\em conformal transformation} is a conformal diffeomorphism from $(M,g)$ to itself.
The conformal transformations form a group and we call a group~$G$ of conformal transformations {\em essential} if there
if there is no metric $\hat g$ in the conformal class of $g$ such that~$G$ is a group of isometries for $\hat g$. Otherwise~$G$ is called {\em inessential}.
A single conformal transformation $\phi$ is called {\em essential} if the group generated by $\phi$ is essential, and inessential otherwise.
A {\em homothetic transformation} or just a {\em homothety} is a conformal transformation for which the function $\vf$ is constant $\vf\equiv\lambda\in \R$.
If $\lambda\not=0$, we say that $\phi$ is a {\em proper homothety}.


% Original source lines 413--666
\section[A Hermitian symmetric space in signature (4n,4n) with compact conformal quotients]{A Hermitian symmetric space in signature $\boldsymbol{(4n,4n)}$ \\ with compact conformal quotients}
 In the following, we use the index conventions
\[A,B,C,\ldots \in \{1,\ldots, 4n\},\qquad a,b,c,\ldots\in \{1\ldots ,2n\},\qquad i,j,k,\ell,\ldots\in \{1,\ldots, n\}.\] We use the Einstein summation convention for these ranges of indices.

Let $(\sigma_{ij})$ be a symmetric complex $n\times n$ matrix and define the quadratic polynomial
\[\sigma\bigl(z^1,\ldots, z^{2n}\bigr):= \sigma_{ij}z^{i}z^{j+n},\]
on $\C^{2n}$.
Using $\sigma$, we define
 $f \colon \C^{4n}\to \rr$ as
\begin{align}
f\bigl(z^1,\ldots , z^{4n}\bigr)
&{}=
\delta_{ab} z^a\ol{z}^{b+2n}+
\sigma\bigl(z^1,\ldots, z^{2n}\bigr)\overline{\sigma\bigl(z^1,\ldots, z^{2n}\bigr)}
\nonumber\\
&{}=
\delta_{ab} z^a\ol{z}^{b+2n}
+
\sigma_{ij}\bar{\sigma}_{k\ell}\bigl(
 z^{i}z^{j+n}\ol{z}^{k} \ol{z}^{\ell+n}\bigr).\label{potential}
\end{align}
As K\"{a}hler potential, this function defines an indefinite K\"ahler metric of signature $(4n,4n)$ on $\C^{4n}=\rr^{8n}$ by
\begin{gather}\label{metric}
g= h_{AB} \,
 \d z^A \cdot \d \ol{z}^B,\qquad \text{where} \quad h_{AB}=\frac{\partial^2 f}{\partial z^A\partial \ol{z}^B},
 \end{gather}
with K\"{a}hler form
 \begin{align*}
\omega
=
\frac{\mathrm{i}}{2} h_{AB}\,
 \d z^A \wedge \d \ol{z}^B.
\end{align*}
The only non-vanishing terms in $h_{AB}$ are the following:
\begin{gather}
h_{a\, b+2n} = h_{a+2n\, b} = \delta_{ab},
\nonumber\\
h_{j\ell} =
\sigma_{ij}\bar{\sigma}_{k\ell}\,
 z^{i+n}\ol{z}^{k+n},
\nonumber\\
h_{j\, \ell+n} =
\sigma_{ij}\bar{\sigma}_{k\ell}\,
 z^{i+n}\ol{z}^{k},
\nonumber\\
h_{j+n\,\ell} =
\sigma_{ij}\bar{\sigma}_{k\ell}\,
 z^{i}\ol{z}^{k+n},
\nonumber\\
h_{j+n\, \ell+n} =
\sigma_{ij}\bar{\sigma}_{k\ell}\,
 z^{i}\ol{z}^{k},
\label{metriccoeff}
\end{gather}
so that {$(h_{AB})$} and its inverse are of the form
\[
{\bigl(h_{AB}\bigr)}=\begin{pmatrix} h_{ab} &\1_{2n} \\ \1_{2n} & 0\end{pmatrix},
\qquad
{\bigl(h^{AB}\bigr)}=\begin{pmatrix} 0 &\1_{2n} \\ \1_{2n} & -h_{ab}\end{pmatrix}.
\]

Without loss of generality, we can assume that the symmetric matrix {$(\sigma_{ij})$} is diagonalised, as the following lemma shows.
\begin{Lemma}\label{diag:lem}
For any complex symmetric $n\times n$ matrix $\sigma=(\sigma_{ij})$, the pseudo-Riemannian manifold $\bigl(\R^{8n},g\bigr)$ is isometric to a pseudo-Riemannian manifold $\bigl(\R^{8n},\hat g\bigr)$ of the same type defined by a diagonal matrix $(\hat\sigma_{ij})$.
\end{Lemma}

\begin{proof}
Assume that $Q\in \GL(n,\C)$ such that $\hat \sigma=Q^\top \sigma Q$ is diagonal. Let $P=\bigl(\overline{Q}^{-1}\bigr)^\top$. The real coordinate transformation $\phi_Q$ induced by complex linear transformation
\[ \hat{z}^{i}:=Q^i{}_k z^{k},\qquad \hat{z}^{i+n}:=Q^i{}_k z^{k+n},\]
and
\[ \hat{z}^{i+2n}:=P^i{}_k z^{k+2n},\qquad {\hat{z}^{i+3n}}:=P^i{}_k z^{k+3n},\]
maps $\R^{8n}$ to $\R^{8n}$ and pulls back the K\"{a}hler potential $f$ to the K\"{a}hler potential $\hat f= f\circ \phi_Q$,
\[
\hat f\bigl({z}^1,\ldots , {z}^{4n}\bigr)
=\delta_{ab} {z}^a\ol{{z}}^{b+2n}+
\hat\sigma_{ij}\bar{\hat\sigma}_{k\ell}
 \bigl(
 {z}^{i}{z}^{j+n}\ol{{z}}^{k} \ol{{z}}^{\ell+n}
 \bigr).
 \]
 As a consequence, the corresponding pseudo-K\"{a}hler metrics are isometric, $\hat g=\phi_Q^*g$. Note also that the quadratic term $\delta_{ab} {z}^a\ol{{z}}^{b+2n}$ is invariant under $\phi_Q$, i.e.,~$\phi_Q\in \O(4n,4n)$.
\end{proof}

We will now compute the Christoffel symbols for the metric (\ref{metric}) defined by the K\"{a}hler potential~(\ref{potential}). From (\ref{metriccoeff}), the metric is given as
\begin{gather}
g =\delta_{ab} \bigl(\d z^a\d \ol{z}^{b+2n} + \d z^{a+2n}\d \ol{z}^{b}\bigr)
\nonumber\\ \hphantom{g=}{}
 +
\sigma_{ij}\bar{\sigma}_{k\ell}
 \bigl(
 z^{i+n}\ol{z}^{k+n}\d z^{j}\d\ol{z}^{\ell}
 +
 z^{i+n}\ol{z}^{k}\d z^{j}\d\ol{z}^{\ell+n}
 +
 z^{i}\ol{z}^{k+n}\d z^{j+n}\d\ol{z}^{\ell}
 +
 z^{i}\ol{z}^{k}\d z^{j+n}\d\ol{z}^{\ell+n}\bigr)
\nonumber\\ \hphantom{g}{}
 =
\delta_{ab} \bigl(\d z^a\d \ol{z}^{b+2n} + \d z^{a+2n}\d \ol{z}^{b}\bigr)
+
 \sigma_{ij} \d\bigl(z^{i} z^{j+n}\bigr)
 \overline{
 \sigma_{k\ell} \d\bigl(z^{k} z^{\ell+n}\bigr)}
.\label{metriclong}
\end{gather}
In the following, we will write $\partial_A$ for $\frac{\partial}{\partial z^A}$ and $\partial_{\overline{A}}$ for $\frac{\partial}{\partial \overline{z}^A}$.
The {(essential)} Christoffel symbols of a K\"{a}hler metric are given by \cite[Section~IX.5]{ko-no2}
\[\Gamma^C_{AB} =h^{DC}\partial_Ah_{BD}.\]
\big(In fact, \smash{$\Gamma^{\bar{C}}_{\bar{A}\bar{B}}= \overline{\bigl(\Gamma^C_{AB}\bigr)}$} and the mixed Christoffel symbols are zero.\big)
In our situation, this shows~that
\[ \Gamma^A_{a+2n\ B}=0,\]
so that the holomorphic vector fields $\del_{a+2n}$, $a=1,\ldots , 2n$, are parallel on $\bigl(\R^{8n}, g\bigr)$. As~a~consequence, the real vector fields
$\del_{a+2n}+\del_{\overline{a+2n}}$ and $\mathrm{i}{\bigl(\del_{a+2n}-\del_{\overline{a+2n}}\bigr)}$ are parallel as well.

Similarly,
one
checks that
\[\Gamma^{c}_{AB} =0.\]
Furthermore, it is easy to see that
\[
\Gamma^{c+2n}_{i j}= \partial_{i}h_{j\, c}=0, \]
as well as
\[
\Gamma^{c+2n}_{i+n\,j+n}=0.
\]
A similar computation shows that the only non-vanishing Christoffel symbols are
\[
\Gamma^{\ell+2n}_{i+n\, j}
 =
 \partial_{i+n} h_{ j\, \ell}
 =
 \partial_{i+n} \bigl(
 \sigma_{mj}\bar{\sigma}_{k\ell}\,
 z^{m+n}\ol{z}^{k+n}\bigr)
 =
 \sigma_{ij}\bar{\sigma}_{k\ell}\,
\ol{z}^{k+n}
\]
and
 \[ \Gamma^{\ell+3n}_{i+n\, j}
 =
 \partial_{i+n} h_{ j\, \ell+n}
 =
 \partial_{i+n} \bigl(
 \sigma_{mj}\bar{\sigma}_{k\ell}\,
 z^{m+n}\ol{z}^{k} \bigr)
 =
 \sigma_{ij}\bar{\sigma}_{k\ell}\,
\ol{z}^{k} .\]
To summarise, for all $b\in \{1, \ldots, 2n\}$ and $A,B\in \{1, \ldots, 4n\}$, we have
\begin{equation}
\label{LC1}
\nabla_{\partial_A}\partial_{b+2n}=0,\qquad
\nabla_{\partial_A}\partial_b \in \Gamma (\mathrm{span} (\partial_{a+2n})_{a=1,\ldots, 2n}).\end{equation}
The curvature tensor is determined by ${\del_{a+2n}\hook R=0}$ and the following equations:
\begin{alignat*}{3}
&R\bigl(\del_{\overline{i}},\del_{j}\bigr)\del_{k} = 0,\qquad&& R\bigl(\del_{\overline{i}},\del_{j}\bigr)\del_{k+n} = \sigma_{jk}\sum_{\ell}\bar{\sigma}_{i\ell} \del_{\ell+3n},&
\\
&R\bigl(\del_{\overline{i+n}},\del_{j+n}\bigr)\del_{k} = \sigma_{jk}\sum_{\ell}\bar{\sigma}_{i\ell} \del_{\ell+2n},\qquad&& R\bigl(\del_{\overline{i+n}},\del_{j+n}\bigr)\del_{k+n} = 0,&
\\
&R\bigl(\del_{\overline{i}},\del_{j+n}\bigr)\del_{k} = \sigma_{jk}\sum_{\ell}\bar{\sigma}_{i\ell} \del_{\ell+3n},\qquad&& R\bigl(\del_{\overline{i}},\del_{j+n}\bigr)\del_{k+n} = 0,&
\\
&R\bigl(\del_{\overline{i+n}},\del_{j}\bigr)\del_{k} = 0,\qquad&& R\bigl(\del_{\overline{i+n}},\del_{j}\bigr)\del_{k+n} =
\sigma_{jk}\sum_{\ell}\bar{\sigma}_{i\ell} \del_{\ell+2n}.&
\end{alignat*}
Using conjugation and the skew symmetry of $R$, we get
\begin{alignat*}{3}
&R\bigl(\del_{\overline{i}},\del_{j}\bigr)\del_{\overline{k}} = 0,\qquad&& R\bigl(\del_{\overline{i}},\del_{j}\bigr)\del_{\overline{k+n}} = -\bar{\sigma}_{ik}\sum_{\ell}{\sigma}_{j\ell} \del_{\overline{\ell+3n}},&
\\
&R\bigl(\del_{\overline{i+n}},\del_{j+n}\bigr)\del_{\overline{k}} = -\bar{\sigma}_{ik}\sum_{\ell}{\sigma}_{j\ell} \del_{\overline{\ell+2n}}, \qquad&& R\bigl(\del_{\overline{i+n}},\del_{j+n}\bigr)\del_{\overline{k+n}} = 0,&
\\
&R\bigl(\del_{\overline{i}},\del_{j+n}\bigr)\del_{\overline{k}} = 0,\qquad&& R\bigl(\del_{\overline{i}},\del_{j+n}\bigr)\del_{\overline{k+n}} =
-\bar{\sigma}_{ik}\sum_{\ell}{\sigma}_{j\ell} \del_{\overline{\ell+2n}},&
\\
&R\bigl(\del_{\overline{i+n}},\del_{j}\bigr)\del_{\overline{k}} =
-\bar{\sigma}_{ik}\sum_{\ell}{\sigma}_{j\ell} \del_{\overline{\ell+3n}},\qquad&& R\bigl(\del_{\overline{i+n}},\del_{j}\bigr)\del_{\overline{k+n}} =0.&
\end{alignat*}
With this information at hand, we obtain the following proposition, for which we recall
that a~semi-Riemannian manifold is indecomposable if it does not split as a product, not even locally.

\begin{Proposition}\label{kaehlerprop}
The metric $g$ in \eqref{metric} defined by the K\"{a}hler potential \eqref{potential} is locally symmetric, K\"{a}hler of neutral signature, Ricci-flat, but not flat and hence not conformally flat. Moreover, the semi-Riemannian manifold $\bigl(\R^{8n},g\bigr)$ is a symmetric space.
\big(More precisely, $\bigl(\C^4=\R^{8n},g\bigr)$ is a Hermitian symmetric space.\big) If $(\sigma_{ij})$ is non-degenerate, then
$\bigl(\R^{8n},g\bigr)$ is indecomposable.
\end{Proposition}

\bprf
Using the fact~\cite[p.~90]{Moroianu07} that
\[ \mathrm{Ric}(\partial_A,\partial_{\bar B}) = -\partial_A\partial_{\bar B}\log \det (g(\partial_C,\partial_{\bar D}))),\]
we see that since the determinant of $h_{AB}$ is equal to $1$, $g$ is Ricci-flat.
With the above formulas~(\ref{LC1}) for the Levi-Civita connection, we have that
\begin{equation*}
%\label{LC}
\nabla_{\partial_A}\partial_B \in \Gamma (\mathrm{span} (\partial_{a+2n})_{a=1,\ldots, 2n}),\qquad \nabla_{\partial_{\bar{A}}}\partial_{\bar{B}} \in \Gamma \bigl(\mathrm{span} \bigl(\partial_{\overline{a+2n}}\bigr)_{a=1,\ldots, 2n}\bigr).\end{equation*}
Together with the fact that the components of the curvature tensor in the coordinates are constant, this implies that the curvature tensor is parallel, so that $g$ is locally symmetric,
$\nabla R=0$.

In order to show that $\R^{8n}$ with the metric $g$ is a globally symmetric space, we have to show that it is geodesically complete.
The geodesic equations are (we only work with the unbarred components, as the barred ones follow by complex conjugation)
\begin{gather*}
0 = \ddot{\gamma}^{a},\qquad\text{ i.e.,}\quad \gamma^a(t)=p^at+q^a,
\\
0 =
 \ddot{\gamma}^{\ell+2n} {+2}
 \sigma_{ij}\bar{\sigma}_{k\ell}\,
\ol{\gamma}^{k+n} \dot\gamma^{i+n}{\dot\gamma^{j}}
=
\ddot{\gamma}^{\ell+2n} {+2}
 \sigma_{ij}\bar{\sigma}_{k\ell}\,
\bigl(p^{k+n} t+q^{k+n}\bigr)
{p^{i+n}p^{j}},
\\
0 =
 \ddot{\gamma}^{\ell+3n} {+2}
 \sigma_{ij}\bar{\sigma}_{k\ell}\,
\ol{\gamma}^{k} \dot\gamma^{i+n}{\dot\gamma^{j}}
=
 \ddot{\gamma}^{\ell+3n} {+2}
 \sigma_{ij}\bar{\sigma}_{k\ell}\,
\bigl(p^{k} t+q^{k}\bigr){p^{i+n}p^{j}}.
\end{gather*}
Hence the $\gamma^{a+2n}(t) $ are cubic polynomials in $t$ and hence
 defined for all $t\in \R$.
As a consequence, $\bigl(\R^{8n},g\bigr)$ is geodesically complete and a symmetric space of signature $(4n,4n)$.

Now we prove that the holonomy algebra of $\bigl(\R^{8n},g\bigr)$ is indecomposable if $(\sigma_{ij})$ is non-degenerate.
Indecomposability means here that there is no decomposition of the (real) tangent space as a proper sum of two complementary non-degenerate invariant subspaces. This is equivalent to the geometric indecomposabilty formulated in the statement of the proposition.

In virtue of Lemma \ref{diag:lem}, we can assume that $(\sigma_{ij})$ is diagonal
with diagonal elements $\lambda_i\neq 0$.
The holonomy algebra is spanned by the following endomorphisms:
\begin{gather*}
R\bigl(\partial_i + \partial_{\bar i},\partial_j+\partial_{\bar j}\bigr) = R\bigl(\partial_{\bar i}, \partial_j\bigr) - R\bigl(\partial_{\bar j}, \partial_i\bigr),\\
R\bigl(\partial_i + \partial_{\bar i},\mathrm{i}\bigl(\partial_j-\partial_{\bar j}\bigr)\bigr) = \mathrm{i}\bigl(R\bigl(\partial_{\bar i}, \partial_j\bigr) +
R\bigl(\partial_{\bar j}, \partial_i\bigr)\bigr),\\
R\bigl(\partial_{i+n} + \partial_{\overline{i+n}},\partial_{j+n}+\partial_{\overline{j+n}}\bigr) = R\bigl(\partial_{\overline{i+n}}, \partial_{j+n}\bigr) - R\bigl(\partial_{\overline{j+n}}, \partial_{i+n}\bigr),\\
R\bigl(\partial_{i+n} + \partial_{\overline{i+n}},\mathrm{i}\bigl(\partial_{j+n}-\partial_{\overline{j+n}}\bigr)\bigr) = \mathrm{i}\bigl(R\bigl(\partial_{\overline{i+n}}, \partial_{j+n}\bigr) + R\bigl(\partial_{\overline{j+n}}, \partial_{i+n}\bigr)\bigr),\\
R\bigl(\partial_{i} + \partial_{\overline{i}},\partial_{j+n}+\partial_{\overline{j+n}}\bigr) = R\bigl(\partial_{\overline{i}}, \partial_{j+n}\bigr) - R\bigl(\partial_{\overline{j+n}}, \partial_{i}\bigr),\\
R\bigl(\partial_{i} + \partial_{\overline{i}},\mathrm{i}\bigl(\partial_{j+n}-\partial_{\overline{j+n}}\bigr)\bigr) = \mathrm{i}\bigl(R\bigl(\partial_{\overline{i}}, \partial_{j+n}\bigr) + R\bigl(\partial_{\overline{j+n}}, \partial_{i}\bigr)\bigr).
\end{gather*}
Inspection of the above formulas for these endomorphisms shows that the holonomy algebra $\mathfrak{hol}$ at the origin is represented
in the basis $(\partial_A)$ of the holomorphic tangent space ${U\!:=\!T_0^{1,0}\C^{4n}\cong \C^{4n}}$ by the algebra of $4n \times 4n$ matrices of the form
\[ \begin{pmatrix}0&0\\\ast &0\end{pmatrix},\]
where $\ast$ stands for an arbitrary skew-Hermitian $2n \times 2n$ matrix. We claim that the real vector space $U$ does not admit any
non-trivial decomposition into complementary subspaces invariant under $\mathfrak{hol}$. Let $U_1\subset U$ be a proper invariant subspace. It is either contained in
\[ U':= \mathrm{span}\{ \partial_{a+2n}\mid a=1,\ldots, 2n\},\]
which is precisely the kernel of $\mathfrak{hol}$, or has a nontrivial projection
to $U/U'$. Considering a line $L:=\R v$, where $v\in U_1 \setminus U'$, we see that in the latter case
\[ U'=\mathfrak{hol}\cdot L.\]
The reason is that the unitary group acts transitively on the unit sphere. By the holonomy invariance of $U_1$, this implies that $U'\subset U_1$, so that any
invariant subspace of $U$ is either contained in $U'$ or contains $U'$. This implies that the intersection
of any two invariant subspaces of complementary dimensions is non-trivial.
\eprf

\setcounter{Theorem}{3}
% Original source lines 689--755
As a consequence of $g$ being locally symmetric, the conformal group of the metric $g$ on $\rr^{8n}$ is contained in the group of homotheties,
\[H=\bigl\{ \phi\in \mathrm{Diff}\bigl(\R^{8n}\bigr)\mid \exists\ \lambda\in \rr\colon \psi^*g=\e^{2\lambda}g\bigr\}.\]
This is due to a result in \cite[Proposition 2.1]{CahenKerbrat77} that states that a conformal transformation between open sets in a semi-Riemannian manifold that is not conformally flat but has parallel Weyl tensor is a homothety. Since $g$ is locally symmetric it must also have parallel Weyl tensor, so that the result follows. Hence, we have the following.

\begin{Proposition}
The conformal group of $\bigl(\R^{8n},g\bigr)$, where $g$ is the metric in \eqref{metric} defined by the K\"{a}hler potential \eqref{potential} is equal to
\[\R\ltimes \Iso\bigl(\R^{8n},g\bigr),
\]
where the $\R$-factor is spanned by a one-parameter group of homotheties $\{\phi_s\mid s\in \R\}$.
Moreover, $g$~admits an abelian $2$-real-parameter family of homothetic transformations induced by the following diagonal linear maps of $\C^{4n}$,
\belabel{chomoth}\phi_{a,b}:=
%\begin{pmatrix} \e^{s+t}\1_{n} &0&0&0\\ 0& \e^{s-t} \1_{n}&0&0 \\
%0&0& \e^{3s-t} \1_{n}&0\\
%0&0&0&\e^{3s+t}\1_{n}
%\end{pmatrix} =
\begin{pmatrix} \e^{a}\1_{n} &0&0&0\\ 0& \e^{b} \1_{n}&0&0 \\
0&0& \e^{a+2b} \1_{n}&0\\
0&0&0&\e^{2a+b}\1_{n}
\end{pmatrix},\qquad a,b\in \R,\end{equation}
that fix the origin $($as any linear map$)$ and satisfy $\phi_{a,b}^*g=\e^{2(a+b)}g$.

\end{Proposition}
\begin{proof}
We begin by showing that the metric $g$ admits a homothety that is not an isometry.
Indeed, for each $s\in \R$, the linear diagonal transformation
\[\phi_s:=\begin{pmatrix} \e^{s}\1_{2n} &0\\ 0& \e^{3s}\1_{2n}\end{pmatrix}\]
of $\C^{4n}$ induces a homothety of $g$ with $\phi_s^*g=\e^{4s}g$.
Composing the homothety $\phi_s$ with the isometry induced by
\[\psi_t:=\begin{pmatrix} \e^{t}\1_{n} &0&0&0\\ 0& \e^{-t} \1_{n}&0&0 \\
0&0& \e^{-t} \1_{n}&0\\
0&0&0&\e^{t}\1_{n}
\end{pmatrix},\]
we obtain the {two-parameter} family of homotheties in the proposition as
\[\phi_{a,b}= \phi_{\tfrac{a+b}{2}} \circ \psi_{\tfrac{a-b}{2}}. \]
This completes the proof.
\end{proof}

We will now construct a compact conformal manifold with essential conformal transformations.

\btheo\label{main:thm}
Let $\tM=\rr^{8n}\setminus\{0\}$ and $g$ the K\"ahler metric on $\tM$ defined by the potential $f$ in~\eqref{potential} for a given non-zero symmetric matrix $(\sigma_{ij})$.
Let $a,b>0$ be fixed positive real numbers and $\Gamma_{a,b} $ be the cyclic group of {holomorphic} homotheties of $(\tM,g)$ generated by $\phi_{a,b}$ in~\eqref{chomoth}. Then
$M_{a,b}=\tM/\Gamma_{a,b}$ is a compact manifold diffeomorphic to $\SS^1\times \SS^{8n-1}$ and carries a $($non-flat$)$ conformal structure of signature $(4n,4n)$ that is locally conformal to a Ricci-flat K\"ahler metric and has essential conformal transformations.
$($The induced integrable complex structure on the quotient manifold preserves the induced conformal structure.$)$
\etheo
\bprf

Since $a$ and $b$ are positive, $M_{a,b}=\tM/\Gamma_{a,b}$ is a compact manifold that is diffeomorphic to $\SS^1\times \SS^{8n-1}$.
In fact, a fundamental domain diffeomorphic to $(0,1)\times \SS^{8n-1}$ is given by
\[ \mathcal D = \Biggl\{ (x_1,x_2,x_3,x_4) \in \mathbb{R}^{8n} = \bigl(\mathbb{R}^{2n}\bigr)^4 \mid 1< \sum_{i=1}^4 \|x_i\|^2 \  \mbox{and}\
\sum_{i=1}^4 \frac{1}{\lambda_i^2} \|x_i\|^2<1\Biggr\},\]
 where $\lambda_1 = \e^a$, $\lambda_2 = \e^b$, $\lambda_3 = \e^{a+2b}$, $\lambda_4 = \e^{2a+b}$, and $\|\, \cdot\,\|$ is the Euclidean norm on $\R^{2n}$.
 Note that this is the region enclosed by the standard sphere and an ellipsoid. Identifying the two
 boundary components of $\mathcal D$ (i.e., the sphere and the ellipsoid) via the diffeomorphism induced by $\phi_{a,b}$ yields a diffeomorphism $M_{a,b} \cong \SS^1\times \SS^{8n-1}$.

{As} $\Gamma_{a,b}$ is a group of homotheties, the metric $g$ defines a conformal structure $\mathbf{c}$ on {$M_{a,b}$.} This conformal structure is locally conformally Ricci-flat K\"ahler and not conformally flat.

For $t\in \rr$, consider the homotheties
\[\phi_{0,t}=
\begin{pmatrix} \1_{n} &0&0&0\\ 0& \e^{t} \1_{n}&0&0 \\
0&0& \e^{2t} \1_{n}&0\\
0&0&0&\e^{t}\1_{n}
\end{pmatrix},
\]
of $\bigl(\tM,g\bigr)$ with fixed point set $\bigl\{\bigl(z^1,\ldots , z^n,0,\ldots, 0\bigr)\mid z^i \in \C^n\setminus \{0\}\bigr\}$. Since the $\phi_{0,t}$ have fixed points, they are essential on $\tM=\rr^{8n}\setminus\{0\}$, see \cite[Proposition 2.5]{LeistnerTeisseire22}.
In addition, for each~${t\in \rr}$, these homotheties commute with $\phi_{a,b}$ and hence descend to a conformal map $\psi_t$ on the quotient~$M_{a,b}$ and are essential on the quotient. Indeed, if there was a metric in the conformal class on~$M_{a,b}$ for which $\psi_t$ {was} an isometry, then this metric would lift to $\tM$ and have~$\phi_{0,t}$ as an isometry. This is a contradiction as the $\phi_{0,t}$ are essential on $\tM$.
\end{proof}
\begin{thebibliography}{99}
\bibitem[3]{CahenKerbrat77}
Cahen M., Kerbrat Y., Transformations conformes des espaces sym\'etriques
 pseudo-riemanniens, \textit{C.~R.~Acad. Sci. Paris S\'er. A-B} \textbf{285}
 (1977), A383--A385.

\bibitem[9]{ko-no2}
Kobayashi S., Nomizu K., Foundations of differential geometry. {V}ol.~{II},
 \textit{Interscience Tracts in Pure and Applied Mathematics}, Vol.~15,
 Interscience Publishers John Wiley \& Sons, Inc., New York, 1969.

\bibitem[10]{LeistnerTeisseire22}
Leistner T., Teisseire S., Conformal transformations of {C}ahen--{W}allach
 spaces, \textit{Ann. Inst. Fourier (Grenoble)}, {t}o appear,
 \href{https://arxiv.org/abs/2201.12958}{arXiv:2201.12958}.

\bibitem[13]{Moroianu07}
Moroianu A., Lectures on {K}\"ahler geometry, \textit{London Math. Soc. Stud.
 Texts}, Vol.~69, \href{https://doi.org/10.1017/CBO9780511618666}{Cambridge University Press}, Cambridge, 2007.

\end{thebibliography}
\end{document}
